\begin{recipe}{Pannenkoeken}
\kicker[Hoofdgerecht,Vegetarisch,Nederlands]{Hoofdgerecht · vegetarisch · Nederlands}
\index[register]{Pannenkoeken}
\index[register]{Bloem}
\index[register]{Eieren}
\index[register]{Melk}

\meta{45 minuten}

\marginimage{images/pannenkoeken}
\lettrine{P}{annenkoeken} eten we niet vaak, vooral als de kinderen op hun
verjaardag mogen kiezen. Om er nog iets gezonds bij te krijgen beginnen we
meestal met een soepje, zoals de simpele kippensoep
(p.~\pageref{recipe:kippensoep-simpel}). Het beslag moet even rusten, dus
begin op tijd.

\blockrule{Ingrediënten}
\begin{ingredients}
  \ing{250 g}{patentbloem}\index[register]{Bloem}
  \ing{500 ml}{melk}
  \ing{2}{eieren}\index[register]{Eieren}
  \ingb{vloeibare boter, voor de pan}
  \ingb{snufje zout}
  \ingb{plakken jong belegen kaas, voor de kaaspannenkoeken (optioneel)}
\end{ingredients}

\blockrule{Bereiding}
\tip{Laat het beslag 20 minuten rusten, dan worden de pannenkoeken soepeler en
scheuren ze minder snel bij het omdraaien.}
\begin{steps}
  \item Doe de bloem en het zout in een kom, maak een kuiltje in het midden
        en breek de eieren erin.
  \item Roer vanuit het midden en voeg de melk beetje bij beetje toe tot een
        glad, dun beslag en laat 20 minuten rusten.
  \item Verhit een koekenpan op middelhoog vuur en smeer de bodem in met een
        beetje vloeibare boter.
  \item Schep een soeplepel beslag in de pan en kantel meteen zodat de bodem
        gelijkmatig bedekt is.
  \item Bak de pannenkoek op middelhoog vuur tot de bovenkant droog ziet en
        de onderkant goudbruin is, draai om en bak nog een minuut. De eerste
        is vaak een proefpannenkoek: zo proef je het beslag en is de pan
        alvast heet en ingevet.
  \item Herhaal voor de rest van het beslag. Van de laatste 5 maken wij vaak
        kaaspannenkoeken: leg na het omdraaien een paar plakken jong belegen
        kaas op de helft en vouw dubbel zodra de kaas smelt. Wij bakken
        vooruit en houden ze warm onder een deksel tot iedereen aan tafel zit.
\end{steps}
\end{recipe}
